\chapter{Análisis de la conexión a Internet}

% ===========================================================
% SECCIÓN 1: INFORMACIÓN GENERAL DEL ROUTER Y CONEXIÓN WAN
% ===========================================================

\section{Descripción General del Router (WAN)}
Esta sección contiene los detalles básicos y las configuraciones físicas de la conexión a Internet.

% Subsección sobre IPs
\subsection*{Direccionamiento IP}
\textbf{1. Dirección IP Pública del Router:} 
\campo{Obtenida a través de la configuración avanzada del router.}
88.24.71.160
\vspace{0.5cm} % Espacio extra para separar grupos de preguntas

\textbf{2. Direcciones Privadas (LAN):} 
% Se pregunta si tiene y luego pide detalles, usando el comando \campo.
\campo{¿Tiene direcciones privadas? Si es afirmativo, especificar la dirección IP privada principal o rango.}
Sí, tiene direcciones privadas. La dirección IP privada principal es 192.168.1.1 y su rango es del 33 al 199

% Subsección sobre filtrado de conexiones
\subsection*{Filtrado de Tráfico (Firewall)}

\textbf{3. ¿El router tiene filtrado de conexiones entrantes?} 
\campo{Respuesta: [Sí / No]} Sí

\vspace{0.2cm} % Pequeño espacio extra

% Uso de un entorno condicional simulado para la explicación
{\bfseries Nota:} En caso afirmativo (y si se especificaron direcciones privadas), responder lo siguiente:

\textbf{4. ¿Hay alguna configuración especial de conexiones entrantes?} 
\campo{Si es afirmativo, detallar las asignaciones específicas (ej.: Puerto Forwarding XYZ).}
Si, tengo un port forwarding configurado para el puerto 3000, ya que tenía un servidor node express montado.
\vspace{0.5cm} % Espacio extra para separar grupos de preguntas

% Subsección sobre contraseña del router
\subsection*{Credenciales y Acceso al Router}

\textbf{5. Origen de la Contraseña del Router:} 
\campo{¿La contraseña asignada es por el instalador o es propia?}
Asignada por el instalador. Ya que es un entorno doméstico no la he cambiado por la falta
de conocimientos de mis familiares.

% ===========================================================
% SECCIÓN 2: DETALLE DE LA CONEXIÓN WI-FI (WLAN)
% ===========================================================

\newpage % Empieza en una nueva página para mejor formato

\section{Configuración de Red Inalámbrica (Wi-Fi / WLAN)}
Detalle sobre la configuración del acceso inalámbrico.

% Subsección de seguridad y medidas adicionales
\subsection*{Seguridad y Funcionalidades Wi-Fi}

\textbf{1. Nivel de Seguridad:} 
\campo{Indicar el nivel utilizado: [WEP / WPA / WPA2 / WPA3]}
WPA2
\vspace{0.2cm} % Pequeño espacio extra

\textbf{2. Medidas Adicionales Implementadas:} 
\campo{Marcar o describir cualquier medida especial (Filtrado de MAC, No anuncio del SSID, desactivación del DHCP/WPS, empleo de 802.X, etc.).}


% Subsección de contraseña Wi-Fi
\subsection*{Credenciales Inalámbricas}

\textbf{3. Origen de la Contraseña WiFi:} 
\campo{¿La contraseña asignada es por el instalador o es propia?}
Asignada por el instalador. Ya que es un entorno doméstico no la he cambiado por la falta
de conocimientos de mis familiares.
\vspace{1cm} % Gran espacio para el cierre del formulario

% ===========================================================
% SECCIÓN 3: ANEXOS (CAPTURA DE PANTALLA)
% ===========================================================

\section*{Anexos y Evidencia Visual}
Adjuntar captura(s) de pantalla de las opciones básicas de configuración del router. Se requiere evidencia visual para validar la información proporcionada en el documento.

\vspace{1cm} 
\hrule % Línea horizontal separadora
\vspace{0.2cm} 
% Puedes añadir un entorno figure o simplemente dejar espacio:
